\section{Introduction}\label{sec:introduction}

Let \(Y\) be a smooth projective complex variety.  Iritani's Gamma integral
structure associates a flat section \(s_Y(E)\) of the quantum connection to a
topological \(K\)-class \(E\) and identifies the flat pairing with the Euler
pairing \cite{IritaniGamma}.  If \(p\in Y\) is a point, then
\[
 [s_Y(E),s_Y(\OO_p))=\chi(E,\OO_p)=\rk(E).              \tag{1.1}
\]
The second entry in (1.1) is therefore a distinguished flat covector.  We
call it the \emph{Gamma point row}.

Birational quantum comparison theorems often split a quantum D-module into
an ambient summand and summands associated with the wall or blow-up centre.
For rank questions, it is unnecessary to identify every entry of the
comparison matrix.  It is enough to know that the point row restricts to the
point row on the ambient summand and vanishes on the supported summands.
This paper separates the exact statements of that form from the additional
analytic input needed to compose several discrepant walls.

The motivating stabilization problem is concrete.  If \(X\) is a smooth
cubic threefold, the case \(m=2\) asks whether
\(X\times\mathbf P^2\) can be rational.  The point-row approach would compare
the primitive-sixth formal-monodromy packet of this product with that of
projective space along a smooth birational factorization.  The one-wall
identities proved below supply the local rank row needed by that comparison.
What is not supplied is a wall-independent sectorial realization at an
intermediate chamber.  We call this the \(X\times\mathbf P^2\) assembly
obstruction.  The terminology names the missing comparison, not a proved
irrationality obstruction for the product.

Our first result concerns a smooth projective simple VGIT wall
\(X_-\dashrightarrow X_+\) in the sense of Gu--Yu--Yu, oriented so that
\(r_+<r_-\) \cite{GuYuYu}.  Their comparison is a formal, pairing-preserving
decomposition
\[
 \QDM(X_-)
 \cong \QDM(X_+)\oplus\bigoplus_j\QDM(S)_j               \tag{1.2}
\]
over an exceptional Laurent completion, where \(S\) is the wall quotient.
Choose corresponding points \(p_\pm\) in the common open set.

\begin{theorem}[Common-open point column]
\label{thm:intro-simple-wall}
At the extremal specialization used in the Gu--Yu--Yu comparison, the
ambient coordinate of the point class is exact when \(r_+<r_-\):
\[
 \pr_{X_+}\Phi([p_-])=[p_+].                              \tag{1.3}
\]
No assertion is made that every wall coordinate of \(\Phi([p_-])\) vanishes.
\end{theorem}

The proof uses the distinguished master-space lift of the common point.  Its
wall restriction vanishes.  Applying the negative-chamber Fourier
isomorphism to the lift and to its wall-frequency derivative shows that the
derivative is zero in the specialized completed source.  Fourier covariance
then makes the positive point image horizontal.  A regular-singular
recursion kills its apparent positive wall-parameter tail.

Our second result is analytic but crepant.  Let \(Y\dashrightarrow Y'\) be a
projective ordinary flop covered by Lee--Lin--Qu--Wang
\cite{LeeLinQuWang}.  Fix one of their analytic-continuation domains in the
extremal variable.

\begin{theorem}[Ordinary-flop point row]
\label{thm:intro-flop}
The analytically continued point section of \(Y\) equals the intrinsic point
section of \(Y'\) on the fixed continuation domain.  Consequently the point
row has the same zero or nonzero restriction to every sectorially realized
formal-monodromy packet whose chosen realization is transported by the graph
gauge on that branch.
\end{theorem}

Pure extremal curves lie in the exceptional locus and cannot meet a point in
the common open set.  This computes the point column exactly at transverse
Novikov degree zero.  A divisor pulled back from the common contraction then
kills the transverse tail coefficient by coefficient.  This is a theorem
about one distinguished column; it does not promote the ancestor comparison
of \cite{LeeLinQuWang} to full descendant invariance.

The discrepant case contains a further frame comparison.  A formal summand,
a sectorial lift of that summand, and the intrinsic large-radius Gamma
section are different objects.  Gu--Yu--Yu prove (1.2) in a Laurent/formal
coefficient ring, while Shen--Shoemaker identify Gamma asymptotic classes at
an extremal specialization \cite{ShenShoemaker}.  Passing from these results
to an intrinsic Gamma point row requires a common pairing-preserving
sectorial realization.  We state the corresponding one-wall consequence
only under that explicit hypothesis.

The negative results show why this distinction is necessary.  A rank-two
incomplete-Gamma connection has a nonzero Stokes shear from a
\(\zetaSix\)-formal line into a rank-visible ambient line.  The example admits
a flat integral hyperbolic doubling and remains after tensoring with
\(\Q[N]/(N^3)\).  Localized partial Fourier--Laplace transform kills
the constant extension which records the shear.  In two Fourier variables,
the delta module at the common origin is invisible after either localization
but transforms back to a constant full-support module.  Boundary support in
Fourier coordinates therefore points in the opposite direction from the
desired rank conclusion.

The surviving two-wall statement is one-sided.  Complete ray-ordered
corrections must have rank-zero \emph{targets}; their sources and scalar
coefficients are irrelevant.  Under that hypothesis, a finite factorization
preserves the point row by elementary linear algebra.  A signed punctual
coefficient of an oriented boundary complex is a plausible shadow of these
factorwise conditions, but it is not proved equivalent to them here.  A
single alternating coefficient can vanish by cancellation; a marked
blockwise comparison is still required.

Sections~\ref{sec:point-row}--\ref{sec:ordinary-flop} prove the positive
results.  Sections~\ref{sec:incomplete-gamma} and
\ref{sec:fourier-boundary} give the countermodels.  Section~\ref{sec:two-wall}
states the conditional assembly theorem, and Section~\ref{sec:scope}
records the exact source and inference boundaries.
